

\chapter{Speculative Designs: NOrec and Beyond}
\section{NOrec: Value‑Based Validation}
\begin{itemize}
    \item Read‑set tracking via timestamps, write‑set buffering.
    \item Validation on every read and at commit time.
    \item PlusCal model and key invariants.
    \item Why `no ownership records': NOrec's write-set holds raw values, validated against a global sequence clock instead of per-location locks.
    \item The read-validate discipline: capture the clock, read the value, re-check the clock --- the same capture$\rightarrow$read$\rightarrow$recheck pattern that prevents partially-visible writes (Part~VII hardware echo).
    \item The invisible-it-committed trap: without value re-reading in validation, a concurrent commit between last read and commit sneaks a stale snapshot through (a real bug class fixed in the companion repo).
    \item Comparison layout: NOrec has the smallest metadata footprint of the in-memory STMs; validation cost grows with read-set size.
    \item Bloom-filter variant (NOrec-BF): a write-set bloom filter lets most reads skip read-set admission entirely --- read-mostly speedup at the cost of rare false positives.
\end{itemize}

\section{Worked Example: The NOrec Read Path}
\index{NOrec}\index{read-validate}\index{read-only transaction}\index{commit order}\index{validation}\index{read-set}
\label{sec:norec-example}

The read-validate discipline in code, mirroring \S\ref{sec:tso-example}'s
capture$\rightarrow$read$\rightarrow$recheck:

\begin{lstlisting}[language=C, caption={NOrec read: capture clock, read, re-check clock.}]
int32_t tm_read_i4(const void *addr) {
    for (;;) {
        uint64_t t1 = atomic_load(&g_clock);       /* capture   */
        int32_t v  = *(const int32_t*)addr;        /* read      */
        uint64_t t2 = atomic_load(&g_clock);       /* re-check  */
        if (t1 == t2) return v;                    /* consistent */
        /* a commit raced us: retry */
    }
}
\end{lstlisting}

\begin{itemize}
    \item Two clock reads bracket the data read; any interleaved commit changes the clock and forces a retry.
    \item This is the read-validate pattern the bullet above names --- the same one Part~II's store-buffer example revealed in hardware.
    \item Note there is no lock and no CAS: the read path is a plain load plus two acquires. That is the entire NOrec point.
    \item Exercise~1 of this chapter invites you to drop the second clock read and watch the invariant tests of Part~III fail.
\end{itemize}

\section{Worked Example: NOrec Bank Transfer}
\index{NOrec}\index{bank transfer}\index{money conservation}\index{account floor}\index{value-based validation}
\label{sec:norec-bank-transfer-example}

The running bank example is deliberately small: transfer money between two
accounts, preserve total money, and reject transfers that would take the
source account below its floor.  In NOrec, transaction-local writes are buffered
until commit; reads first consult the write-set, then perform the
capture$\rightarrow$read$\rightarrow$recheck protocol.

\begin{lstlisting}[language=C, caption={A compilable C++17 NOrec-style bank transfer with value revalidation.}]
#include <atomic>
#include <cstdint>
#include <mutex>
#include <unordered_map>
#include <vector>

struct ReadEntry {
    const int64_t *addr;
    int64_t value;
};

struct WriteEntry {
    int64_t *addr;
    int64_t value;
};

struct Tx {
    uint64_t start_clock = 0;
    std::vector<ReadEntry> read_set;
    std::vector<WriteEntry> write_set;
};

static std::atomic<uint64_t> g_clock{0};
static std::mutex g_commit_lock;

static bool validate(const Tx &tx) {
    for (const ReadEntry &r : tx.read_set) {
        if (*r.addr != r.value) return false;
    }
    return true;
}

static bool find_write(const Tx &tx, const int64_t *addr, int64_t &out) {
    for (auto it = tx.write_set.rbegin(); it != tx.write_set.rend(); ++it) {
        if (it->addr == addr) {
            out = it->value;
            return true;
        }
    }
    return false;
}

static int64_t tx_read(Tx &tx, int64_t *addr) {
    int64_t buffered = 0;
    if (find_write(tx, addr, buffered)) return buffered;

    for (;;) {
        uint64_t c1 = g_clock.load(std::memory_order_acquire);
        int64_t v = *addr;
        uint64_t c2 = g_clock.load(std::memory_order_acquire);

        if (c1 == c2 && validate(tx)) {
            tx.read_set.push_back(ReadEntry{addr, v});
            return v;
        }

        tx.read_set.clear();
        tx.start_clock = g_clock.load(std::memory_order_acquire);
    }
}

static void tx_write(Tx &tx, int64_t *addr, int64_t value) {
    for (WriteEntry &w : tx.write_set) {
        if (w.addr == addr) {
            w.value = value;
            return;
        }
    }
    tx.write_set.push_back(WriteEntry{addr, value});
}

static bool tx_commit(Tx &tx) {
    std::lock_guard<std::mutex> guard(g_commit_lock);

    if (!validate(tx)) return false;

    for (const WriteEntry &w : tx.write_set) {
        *w.addr = w.value;
    }

    g_clock.fetch_add(1, std::memory_order_release);
    return true;
}

static bool transfer(int64_t accounts[2],
                     int64_t floor[2],
                     int from,
                     int to,
                     int64_t amount) {
    Tx tx;
    tx.start_clock = g_clock.load(std::memory_order_acquire);

    int64_t src = tx_read(tx, &accounts[from]);
    int64_t dst = tx_read(tx, &accounts[to]);

    if (src - amount < floor[from]) return false;

    tx_write(tx, &accounts[from], src - amount);
    tx_write(tx, &accounts[to], dst + amount);

    return tx_commit(tx);
}

int main() {
    int64_t accounts[2] = {100, 100};
    int64_t floor[2] = {0, 0};

    bool ok = transfer(accounts, floor, 0, 1, 25);

    if (!ok) return 1;
    if (accounts[0] + accounts[1] != 200) return 2;
    if (accounts[0] < floor[0] || accounts[1] < floor[1]) return 3;

    return 0;
}
\end{lstlisting}

\begin{itemize}
    \item The money-conservation invariant is checked after commit: the sum of
    committed account balances remains unchanged.
    \item The floor check is speculative.  It is safe only because commit
    revalidates the values that justified it.
    \item The read path validates the old read-set before admitting a new read.
    Otherwise a transaction can assemble a snapshot from different commit
    times.
    \item NOrec has no per-account ownership record; the commit lock and global
    sequence clock serialize publication.
\end{itemize}

\begin{figure}[htbp]
\centering
\begin{tikzpicture}[
    font=\small,
    node distance=1.8cm,
    box/.style={draw, rectangle, rounded corners, minimum width=2.6cm, minimum height=0.8cm, align=center},
    edge/.style={-{Stealth}, thick}
]
\node[box] (begin) {Begin\\read clock};
\node[box, right=of begin] (read) {Read account\\record value};
\node[box, right=of read] (validate) {Validate\\read-set};
\node[box, right=of validate] (buffer) {Buffer\\debit/credit};
\node[box, below=of validate] (commit) {Commit lock\\revalidate};
\node[box, left=of commit] (abort) {Abort\\retry};
\node[box, right=of commit] (publish) {Write back\\advance clock};

\draw[edge] (begin) -- (read);
\draw[edge] (read) -- (validate);
\draw[edge] (validate) -- node[above] {ok} (buffer);
\draw[edge] (buffer) -- (commit);
\draw[edge] (commit) -- node[above] {ok} (publish);
\draw[edge] (validate) -- node[left] {changed} (abort);
\draw[edge] (commit) -- node[above] {changed} (abort);
\draw[edge] (abort.west) -| ($(begin.south west)+(-0.4,-0.2)$) -- (begin.south);
\end{tikzpicture}
\caption{NOrec bank-transfer state machine.  The transaction may validate during reads and must validate again while holding the commit lock; only then may it publish both account updates and advance the global clock.}
\label{fig:norec-bank-state-machine}
\end{figure}


\section{TL2: The Textbook Global-Clock STM}
\begin{itemize}
    \item Global version clock read at begin; snapshot = transaction's start time.
    \item Write-set of (address, value) with per-location version metadata; read-set of (address, observed-value).
    \item Commit: acquire commit lock, re-validate read-set against the clock, write back under lock, advance clock, release.
    \item Timeline visualization (figure): two concurrent TL2 transactions and the clock values they observe.

\begin{figure}[ht]
\centering
\begin{tikzpicture}[font=\small]
  % global clock timeline
  \draw[-{Stealth}] (0,0) -- (11,0) node[right] {global clock};
  \foreach \x/\lab in {1/$5$,3/$5$,5/$6$,7/$6$,9/$7$}
    \draw (\x,0.08) -- (\x,-0.08) node[below] {\lab};
  % T1: begins at 5, reads, then aborts after clock advances
  \node[draw,rectangle,minimum height=0.5cm,minimum width=3.2cm,fill=black!6]
        (T1) at (2.1,1.0) {\texttt{read x}};
  \draw[-{Stealth}] (T1.south) -- (1.4,0.05);
  \draw[-,thick,black!30] (1.4,0.05) -- (5.4,0.05);
  \node at (4.6,1.0) {$T_1$ begin@5};
  % T2: begins at 5, writes x, commits at clock 6
  \node[draw,rectangle,minimum height=0.5cm,minimum width=3.2cm,fill=black!6]
        (T2) at (4.4,2.2) {\texttt{write x}};
  \draw[-{Stealth}] (T2.south) -- (3.4,0.05);
  \draw[-,thick,black!50] (3.4,0.05) -- (5.0,0.05);
  \draw[-{Stealth},thick] (5.0,0.05) -- (5.4,0.05);
  \node at (6.6,2.2) {$T_2$ begin@5, commit@6};
  \node[align=left,text width=7.2cm] at (7.6,1.0)
        {$T_1$ re-validates at commit: version of \texttt{x}\\
         moved past its snapshot $\Rightarrow$ abort};
\end{tikzpicture}
\caption{TL2 timeline: $T_1$ and $T_2$ both begin at clock 5; $T_2$ commits (advancing the clock to 6) before $T_1$ validates, so $T_1$ must abort.}
\label{fig:tl2-timeline}
\end{figure}

    \item Validation semantics: every read-set entry must still match its recorded version, proving no writer committed during the transaction --- this is what makes invisible reads safe.
    \item Why TL2 is the standard baseline: minimal algorithm, clear invariants, easy to model-check (PlusCal model in the companion repo).
    \item The stored-value trap: without value-based re-validation, clock-only validation misfires on OCC write-back races --- the read-validate pattern recurs (see ROMULUS/MVLog case studies).
\end{itemize}

\section{TinySTM Design Variants}
\begin{itemize}
    \item Write‑through vs. write‑back (lazy version management).
    \item Eager locking vs. commit‑time locking.
    \item Comparison of strategies under different workloads.
    \item WBCTL (write-back, commit-time locking) vs.\ WT (write-through): the observable difference is \emph{where} the data lives before commit.
    \item Rollback mechanics: undo log for write-through, discard of write-set for write-back.
    \item Optimizations for the cache: versioned write-set entries, `read-set bypass' when a value matches the write-set, and when each pays off.
\end{itemize}

\section{Multiversion STMs: JVSTM and MVLog}
\begin{itemize}
    \item JVSTM: versioned boxes (VBox) store a history of (version, value) pairs; readers walk the newest body with version $\le$ read-version.
    \item Read-only transactions never abort: a consistent snapshot is read from history rather than validated.
    \item Commit: global clock increments under a lock; writer appends a new VBox body; validation compares VBox heads.
    \item MVLog: commit-slot negotiated at begin (a single \texttt{fetch\_add}), an ever-growing commit log of write-sets, and per-address newest-writer tracking.
    \item MVLog's `no-false-negative' bloom filter lets reads skip the log when the address is provably clean.
    \item The cost of history: memory grows with transaction count; reclamation (epochs, hazardous pointers) becomes a first-order concern.
    \item Why multiversion dominates GPU TM: read-only transactions never abort, and commit slots are cheap atomics (Part~VII, CSMV and GUST).
\end{itemize}

\section{Comparing the Family}
\index{comparison}\index{version management}\index{conflict detection}
\begin{itemize}
    \item A single table mapping design axes --- conflict detection (eager vs.\ lazy), version management (write-through vs.\ write-back), metadata (per-location locks vs.\ global clock), read set (invisible vs.\ validation-driven).
    \item NOrec bread-and-butter simplicity; TL2 clock-heavy but robust; TinySTM tuning knobs; multiversion for read-heavy and GPU workloads.
    \item The recurring invariant: every design must prevent a snapshot from becoming inconsistent mid-flight; the read-validate pattern is how software enforces it.
\end{itemize}

\begin{table}[ht]
\centering
\small
\begin{tabular}{lp{3.1cm}p{3.4cm}p{3.4cm}}
\hline
\textbf{Design} & \textbf{Read-set validation} & \textbf{Write-set ownership} & \textbf{Commit linearization} \\
\hline
SGL & none (mutex)           & direct writes under lock      & lock release \\
NOrec & clock-capture re-read & buffered, no locks            & clock advance \\
TL2 & per-location version   & buffered, commit lock         & clock advance + write-back \\
TinySTM WBCTL & lock/version check & commit-time locking           & lock release \\
TinySTM WT & versioned metadata   & write-through + undo log       & lock release \\
JVSTM & none (history walk)   & VBox append under lock        & clock increment \\
MVLog & no-false-negative bloom & commit-log slot + index      & slot publish \\
\hline
\end{tabular}
\caption{The comparison matrix used throughout this chapter and the exercises: each row gives the same four facts --- read validation rule, write ownership, and the linearization point. A backend ``family'' is a choice in each column; opacity arguments in Part~III are written column-by-column.}
\label{tab:family-matrix}
\end{table}

\begin{itemize}
    \item \textbf{Per-runtime performance signatures} (the expected bottleneck of each row in Table~\ref{tab:family-matrix}): SGL serializes everything (throughput cap $=1/\text{tx time}$); NOrec validates whole read-sets (cost grows with read-set size); TL2 pays commit-lock serialization plus clock reads; TinySTM pays per-location metadata traffic; multiversion pays memory reclamation.
    \item Exercise~6 turns Table~\ref{tab:family-matrix} into numbers: the memory-overhead row is measurable on the repository's benchmarks.
\end{itemize}

\section{Backend Taxonomy: The Repository in One Table}
\index{backend taxonomy}\index{SGL}\index{TinySTM}\index{WBCTL}\index{NOrec}\index{TL2}\index{MVLog}\index{JVSTM}\index{SwissTM}\index{Romulus}\index{XTM}\index{TSX}\index{GPU TM}\index{persistent TM}\index{distributed TM}\index{reclamation}\index{private memory}\index{memory reclamation}
The companion repository contains 20+ runtimes. This book explains the
families; the table maps \emph{which repository backend is which family},
so a reader looking for ``where is the code for X'' can find it and so the
prose never overclaims that every backend is analysed in depth:

\begin{table}[ht]
\centering
\small
\begin{tabular}{lll}
\hline
\textbf{Repository backend} & \textbf{Family} & \textbf{Chapter} \\
\hline
\texttt{sgl}                 & single global lock   & Ch.~8 \\
\texttt{tsxsgl}, \texttt{spht} & HTM + SGL fallback  & Ch.~8, 16 \\
\texttt{norec}, \texttt{norec\_bf} & value-based validation & Ch.~9 \\
\texttt{tl2}, \texttt{tsc\_tm}     & global-clock OCC     & Ch.~9 \\
\texttt{tinystm} (\texttt{wbctl}/\texttt{wt}/\texttt{wbetl}) & per-object locking & Ch.~9 \\
\texttt{jvstm}, \texttt{mvlog}      & multiversion         & Ch.~9 \\
\texttt{swisstm}, \texttt{romulus}, \texttt{leftright} & OCC variants & Ch.~9 \\
\texttt{xtm}, \texttt{dudetm}, \texttt{nvhtm} & page/persistent/hybrid & Ch.~13 \\
\texttt{persistent\_sgl}, \texttt{distributed\_sgl}, \texttt{tikv} & persistent/distributed & Ch.~18 \\
\texttt{gust}, \texttt{csmv}, \texttt{jvstm-gpu}, \texttt{gputx}, \texttt{gacco}, \texttt{gpu\_stm} & GPU warp TM & Ch.~17 \\
\hline
\end{tabular}
\caption{Repository backend $\rightarrow$ family $\rightarrow$ where it is discussed. Backends without a dedicated prose section are still exercised by the benchmark suite and, where present, have a PlusCal/TLA+ model (Table~\ref{tab:traceability}).}
\label{tab:backend-taxonomy}
\end{table}

\section{Semantic Obligations Checklist}
\label{sec:semantic-obligations}
Part~I introduced opacity, sandboxing, and privatization; here is where
each is \emph{discharged} by a concrete runtime mechanism, so no promise
from Chapter~2 floats unresolved:

\begin{itemize}
    \item \textbf{Opacity} $\rightarrow$ Ch.~9 read-validate discipline (\S\ref{sec:norec-example}) and the TLA+ invariants of Table~\ref{tab:traceability}.
    \item \textbf{Sandboxing} $\rightarrow$ Ch.~13 memory management: speculative writes must live in the write-set or the TM region, never the caller's heap.
    \item \textbf{Privatization} $\rightarrow$ Ch.~9 (publishing objects) and Ch.~13: a committed object handed to a non-transactional thread must be fully written back before publication.
    \item \textbf{Irrevocability} (operations that cannot abort, e.g.\ I/O) $\rightarrow$ Ch.~8's serializing designs and Ch.~13's policy.
    \item \textbf{Memory reclamation} (when may a freed address be reused) $\rightarrow$ Ch.~13: deferred frees and the double-free hazards of region allocators.
    \item \textbf{Signal safety} (what \texttt{siglongjmp} may do inside a transaction) $\rightarrow$ Ch.~13's policy section.
    \item \textbf{Compiler obligations} (which accesses must stay instrumented) $\rightarrow$ the contract of Chapter~11.
\end{itemize}

\section{Worked Example: The TL2 Commit Path}
\index{TL2}\index{commit protocol}\index{commit lock}\index{global clock}\index{OCC}\index{validation}\index{version management}
\label{sec:tl2-example}

The five-phase commit that Chapter~6's model and Part~III's invariants check:

\begin{lstlisting}[language=C, caption={TL2 commit: lock, validate, write back, publish.}]
bool tm_end(void) {
    if (!commit_lock.try_lock()) return retry();        /* 1. serialise */
    for (each r in read_set)                            /* 2. validate */
        if (version_of(r.addr) != r.observed_version)
            { commit_lock.unlock(); return retry(); }
    uint64_t ts = global_clock.fetch_add(1) + 1;        /* 3. new version */
    for (each w in write_set) {                         /* 4. write back */
        *(w.addr) = w.value;
        store_version(w.addr, ts);
    }
    commit_lock.unlock();                               /* 5. publish   */
    return true;
}
\end{lstlisting}

\begin{itemize}
    \item Step 2 is the expensive one: the read-set is validated against per-location version metadata, not a global clock.
    \item Step 4 writes back \emph{under} the commit lock, then advances the clock so no concurrent reader can observe a half-published write-set.
    \item The stored-value trap is step 4 interacting with step 2: without the read-validate pattern of \S\ref{sec:norec-example}, a stale read-set can validate against an old version while the write-back overwrites a newer commit.
    \item The same five steps appear in the PlusCal model of Part~III, and in the GPU backends of Part~VII at warp granularity.
\end{itemize}

\section{Worked Example: Account-Floor PlusCal Model}
\index{PlusCal}\index{TLA+}\index{model checking}\index{account floor}\index{money conservation}
\label{sec:pluscal-bank-floor-example}

The model below is intentionally small enough for TLC.  Two processes repeatedly
try to transfer one unit in opposite directions.  The commit action is atomic:
it checks the floor against the current state and then updates both accounts.
This is the serial specification that NOrec, TL2, TinySTM, JVSTM, and MVLog
must refine.

\begin{lstlisting}[language={}, caption={PlusCal/TLA+ model of two-account transfer with conservation and floor invariants.}]
---- MODULE BankTransferFloor ----
EXTENDS Naturals, TLC

CONSTANT Init0, Init1, Floor0, Floor1

ASSUME Init0 >= Floor0 /\ Init1 >= Floor1

(*
--algorithm Bank
variables bal = [0 |-> Init0, 1 |-> Init1];

define
  TotalOK == bal[0] + bal[1] = Init0 + Init1
  FloorsOK == bal[0] >= Floor0 /\ bal[1] >= Floor1
end define;

process (P \in {0, 1})
variables from = P, to = 1 - P, amount = 1;
begin
Loop:
  while TRUE do
    if bal[from] - amount >=
       IF from = 0 THEN Floor0 ELSE Floor1 then
      bal[from] := bal[from] - amount;
      bal[to] := bal[to] + amount;
    end if;
  end while;
end process;
end algorithm;
*)

==== 
\end{lstlisting}

A minimal TLC configuration is:

\begin{lstlisting}[language={}, caption={TLC configuration for the PlusCal account-floor model.}]
SPECIFICATION Spec
CONSTANTS
  Init0 = 2
  Init1 = 2
  Floor0 = 0
  Floor1 = 0
INVARIANTS
  TotalOK
  FloorsOK
\end{lstlisting}

\begin{itemize}
    \item The invariant \texttt{TotalOK} is the money-conservation property:
    transfers move value; they do not create or destroy it.
    \item The invariant \texttt{FloorsOK} is the safety property enforced by
    the transactional read of the source account followed by validation.
    \item To model a broken STM, split the guarded update into separate read,
    validate, debit, and credit actions, then remove the final validation.
    TLC should find an execution where the guard was true for a stale source
    value.
    \item This model is a serial specification, not an implementation model.
    The implementation models add clocks, read-sets, write-sets, locks, or
    version histories.
\end{itemize}

\begin{table}[htbp]
\centering
\small
\begin{tabular}{llll}
\hline
\textbf{Backend family} & \textbf{Snapshot mechanism} & \textbf{Read-only aborts} & \textbf{Main cost} \\
\hline
NOrec & global clock + value re-read & possible & read-set validation \\
NOrec-BF & global clock + bloom-filter bypass & possible & false positives and validation \\
TL2/TSC-TM & start timestamp + per-location versions & possible & metadata traffic and commit lock \\
TinySTM WBCTL & per-location lock/version words & possible & ownership metadata \\
TinySTM WT & write-through + undo log & possible & rollback and eager conflicts \\
JVSTM & version history walk & no & history retention \\
MVLog & commit log + newest-writer index & no for stable snapshots & log scan and reclamation \\
\hline
\end{tabular}
\caption{Snapshot mechanisms in the speculative STM family.  The measured repository results match the expected signatures: TL2 and TSC-TM are close on the M1 Pro bank/fuzz workload, NOrec-BF improves read-mostly NOrec, and MVLog pays for history and log management.}
\label{tab:snapshot-taxonomy}
\end{table}


\section{Summary}
\begin{itemize}
    \item Speculation unlocks concurrency but introduces aborts.
    \item Different version‑management policies trade off simplicity for performance.
    \item The design space is small: clocks, locks, logs, and history --- every real STM is a combination of these.
\end{itemize}

\section{Exercises}
\begin{enumerate}
    \item In NOrec, T1 writes to \texttt{x}, then T2 writes to \texttt{x} and commits. Explain why T1 will abort.
    \item TinySTM can be configured for eager locking or commit‑time locking. Given a workload with high contention on a single location, predict which mode will perform better and justify your answer.
    \item Design a simple experiment using the STAMP benchmark suite to compare the abort rates of NOrec and TinySTM’s write‑through variant.
    \item In TL2, explain why the read-set validation must re-read the metadata (not just remember the snapshot) and sketch the read-validate pattern that prevents a partially committed write from being observed.
    \item Write a PlusCal model of a TL2 commit that omits the read-set re-validation, and run TLC to produce a counterexample (Part~III workflow).
    \item Compare memory overhead: NOrec, TL2, JVSTM, and MVLog, for a workload with a large read-set and small write-set.
    \item \emph{[implementation]} Modify the NOrec bank-transfer example in \S\ref{sec:norec-bank-transfer-example} to retry aborted transactions in a bounded loop.  Report separately the number of floor rejections and validation aborts.
    \item \emph{[verification]} Extend the PlusCal model in \S\ref{sec:pluscal-bank-floor-example} with a speculative local variable \texttt{seen}.  Split the transfer into read, check, debit, and credit actions, then add a commit-time validation action.  Show which invariant fails when the validation action is removed.
    \item \emph{[theory]} In Table~\ref{tab:snapshot-taxonomy}, identify which families implement invisible reads.  For each such family, state the proof obligation that prevents an inconsistent snapshot from being committed.
\end{enumerate}

%  PART V: COMPILER INSTRUMENTATION FOR TM
